\section{Uniformity over coefficient fields}\label{sec:uniformity}

The pointwise rank limit does not by itself allow the field or
coefficients to change at each level. We prove the required uniformity
in two steps. This section first works with any group $\Gamma$ and
compatible finite quotients $Q_i$, provided every refinement kernel
is a $p$-group.

\begin{lemma}[$p$-refinement monotonicity]\label{lem:monotone}
\notformalized
Let $K$ have characteristic $p$, and fix a square matrix
$M\in\operatorname{Mat}_B(K[\Gamma])$ with $B>0$.
If $Q_j\twoheadrightarrow Q_i$ has kernel $P$, a finite $p$-group,
then
\[
 \dim_K\ker M_{Q_j}\le |P|\dim_K\ker M_{Q_i},\qquad
 \frac{\dim_K\ker M_{Q_j}}{Bq_j}
       \le\frac{\dim_K\ker M_{Q_i}}{Bq_i}.
\]
The field and coefficients are fixed within this comparison.
\end{lemma}
\begin{informalproof}
Choose a section $s:Q_i\to Q_j$. Each fine sheet is uniquely $s(x)k$
with $k\in P$, and left multiplication by $g$ takes it to
\[
 s(g_ix)\kappa(g,x)k,\qquad
 \kappa(g,x)=s(g_ix)^{-1}g_js(x)\in P.
\]
If $M_{ab}=\sum_gc_{ab,g}g$, its fine block from coarse column $(b,x)$
to coarse row $(a,y)$ is the operator
\[
 \sum_{g:g_ix=y}c_{ab,g}\rho_P(\kappa(g,x))\quad\text{on }K[P].
\]

Every nonzero finite-dimensional $K[P]$-module has a nonzero fixed
vector. Induct on $|P|$: a nontrivial finite $p$-group has a central
element $z$ of order $p$ (by its class equation and Cauchy's theorem),
and $(z-I)^p=0$. Its nonzero kernel is $P$-invariant, with action
factoring through $P/\langle z\rangle$; induction applies there.
For trivial $P$ there is nothing to prove. Repeating on quotients
gives a full flag in $K[P]$ with one-dimensional trivial factors.

Tensor this flag with the coarse coordinate space. On every quotient
layer, each $\rho_P(\kappa)$ acts as $1$, so the induced matrix is
exactly $M_{Q_i}$. All colliding words contribute their scalar sum;
none is discarded. Apply \eqref{eq:filtered-kernel} to the $|P|$
layers and use $q_j=|P|q_i$.
\end{informalproof}

\begin{theorem}[Finite-type uniformization]\label{thm:uniformization}
\notformalized
Let $\mathscr R$ be a finite-type commutative unital $\F_p$-algebra,
and let $M\in\operatorname{Mat}_B(\mathscr R[\Gamma])$, $B>0$, have
fixed finite support. Use compatible finite quotients with $p$-group
refinement kernels. For a unital map $a:\mathscr R\to K$ to a field,
write
\[
 r_i(a)=\frac{\dim_K\ker M_{Q_i}(a)}{Bq_i}.
\]
Suppose $0\le\beta<1$ and $\limsup_i r_i(a)\le\beta$ for every
such map to a \emph{finite} field. For every $\varepsilon>0$ there
is $I$ such that
\[
 r_i(a)<\beta+\varepsilon
\]
for all $i\ge I$, all characteristic-$p$ fields $K$, and all unital
maps $a:\mathscr R\to K$.
\end{theorem}
\begin{informalproof}
The zero ring has no unital maps to fields. If
$\theta=\beta+\varepsilon>1$, the assertion follows from $r_i\le1$.
Otherwise $0<\theta\le1$. Set $N_i=Bq_i$,
$k_i=\lfloor(1-\theta)N_i\rfloor$, and let $J_i$ be the ideal of
$\mathscr R$ generated by all $(k_i+1)$-minors of $M_{Q_i}$.
For every field specialization,
\begin{equation}\label{eq:bad-minors}
 r_i(a)\ge\theta\quad\Longleftrightarrow\quad
 \rank M_{Q_i}(a)\le k_i
 \quad\Longleftrightarrow\quad a(J_i)=0.
\end{equation}
The rounding is exact, including $\theta=1$.

If $J=\sum_iJ_i$ were proper, choose a maximal ideal
$\mathfrak m\supseteq J$. By the finite-type weak
Nullstellensatz~\cite[Theorem~10.34.1]{Stacks00FV},
$\mathscr R/\mathfrak m$ is a finite extension of $\F_p$.
Its residue map would be bad at every level by \eqref{eq:bad-minors},
contradicting the pointwise hypothesis. Thus $J=\mathscr R$ and
there is a finite identity
\begin{equation}\label{eq:unit-identity}
 1=\sum_{\ell=1}^m f_\ell\Delta_\ell,
\end{equation}
where $\Delta_\ell$ is a generating minor at some level $i_\ell$.
Take $I=\max_\ell i_\ell$.

If a specialization were bad at $j\ge I$, then
\cref{lem:monotone}, applied to that same field and tuple, would
make it bad at every $i_\ell$. All $a(\Delta_\ell)$ would vanish,
contradicting \eqref{eq:unit-identity}. This proves the claim.
The ideals $J_i$ need not be nested; only their bad field loci are
ordered by monotonicity.
\end{informalproof}

This argument also allows equations among parameters and finitely
many nonzero constraints, represented by inverse variables. It supplies
an eventual threshold, not a useful numerical estimate for it. A
search for \eqref{eq:unit-identity} is possible when the ring, chain,
and matrices are effectively presented and the threshold is rational;
no running-time bound is asserted.

\begin{theorem}[Field-uniform half rate]\label{thm:fixed-width}
\notformalized
Fix the template, embedding, and chain of \cref{sec:model}, and
suppose \eqref{eq:separation}. Fix widths $w_v\ge0$ of total $B>0$.
For every $\varepsilon>0$ there is $I$ such that for every $i\ge I$,
every characteristic-$p$ field $K$, and every homogeneous coefficient
tuple fitting the template,
\[
 \frac{\dim_K\ker M_{Q_i}}{Bq_i}<\frac12+\varepsilon.
\]
The threshold is independent of $K$ and the coefficient tuple.
\end{theorem}
\begin{informalproof}
Give each scalar entry of each $C_e$ an independent variable and
form the universal check over their polynomial ring over $\F_p$.
Every finite-field tuple has pointwise rate at most $1/2$ by
\cref{thm:prime-field} and restriction of scalars. Remove zero-width
types if necessary; the remaining incidence graph retains separation.
Apply \cref{thm:uniformization} with $\beta=1/2$.
\end{informalproof}

This rules out above-half homogeneous constructions with fixed width
even when they change fields and coefficients at every size. It does
not bound arbitrary sheet-dependent maps or nonlinear code supports.
The $p$-kernel condition is material: in characteristic three, $1+a$
has nullity zero on the trivial quotient but positive nullity on $C_2$.
Thus \cref{lem:monotone} is false for unrestricted refinements.
